% ============================================================================
%  摘要与关键词
% ============================================================================
\section*{摘要}
\addcontentsline{toc}{section}{摘要}
\markboth{摘要}{摘要}

本文围绕实 Grassmann 流形的欧氏非嵌入问题，系统整理典范丛、正交补丛、切丛公式与
Stiefel--Whitney 特征类之间的关系。对 $\GK=G_k(\R^{k+n})$，首先通过局部图像坐标
复现切丛同构 $T\GK\cong\Hom(\cG_k,\cGc_k)$；再由法丛的稳定平凡化说明：若闭流形
$M^d$ 嵌入 $\R^{d+c}$，则法丛的总 Stiefel--Whitney 类等于切丛总类的逆。由此，
非零的共轭类 $\wbar_j(TM)$ 为余维提供必要的下界；在闭流形的欧氏嵌入临界次数判据
下，它还排除余维恰等于 $j$ 的端点。本文特别强调：形式多项式的计算
必须在真实的 $H^*(\GK;\Ftwo)$ 商环中约化，且正交补丛的 Stiefel--Whitney 类、
切丛的共轭类、以及嵌入与浸入这三组概念不可混用。

在此基础上，本文核验并复现 Oproiu 关于 $G_{2,n}$ 与 $G_{3,n}$ 的代表性非嵌入
论证，说明二进制参数分段如何与 Schubert 类的非零性、以及欧氏嵌入的临界次数条件
相结合；同时整理 Hiller--Stong 关于实 Grassmann 流形浸入维数的命题与界，并将其
严格保持在 immersion 的语境之内。最后，通过 $G_{2,2}$ 等小参数例子展示
“总类公式——商环约化——非零性判定——环境维数排除”这一计算链条，并区分已知定理、
个人复现与尚待证明的探索性观察。本文的结论是：Stiefel--Whitney 类能够提供可靠的
非嵌入必要条件与浸入维数下界，但通常不足以单独决定最小的欧氏嵌入维数。

\vspace{0.8em}
\noindent\tb{关键词：}实 Grassmann 流形；典范丛；切丛；Stiefel--Whitney 类；
共轭特征类；Schubert 类；Gr\"obner 基；欧氏嵌入；浸入；法丛

\vspace{1.2em}
\begin{scope}
本文是《文献阅读与课题调研》课程的持续写作主稿，其目标是形成一份结构完整、
数学论证严谨、引用可核验、可提交的个人读书报告。每周的文献阅读属于研究准备，
其成果须写入本文，而不应只留下独立的阅读笔记。

写作遵循两条原则：其一，先形成可读的初稿，再随文献核验持续修订；其二，明确区分
已核验内容、待核验内容与个人探索。全文采用统一记号 $\GK=G_k(\R^{k+n})$，其维数
为 $kn$；若所引文献采用不同记号，引用时均明确给出转换关系。
\end{scope}

\clearpage
